\chapter{El fin de la esclavitud operativa}

Imagina una aldea antigua. Cada día, sus habitantes deben caminar kilómetros hasta el río, llenar vasijas de barro y cargar el agua de regreso bajo el sol. Es un trabajo agotador, repetitivo y, sobre todo, esencial. Si el aldeano no carga la vasija, no hay agua. Este es el estado actual del desarrollo de software: el programador como el aldeano, cargando cada línea de código, revisando cada error, moviendo datos de un lado a otro manualmente.

Luego llegaron los ingenieros romanos. Ellos no se preguntaron cómo hacer que el aldeano cargara la vasija más rápido o cómo hacer vasijas más ligeras. Se preguntaron: ¿cómo podemos diseñar un entorno donde el agua llegue sola?

El acueducto es la victoria de la arquitectura sobre la ejecución manual. Una vez construido, el acueducto no requiere que nadie cargue nada. El sistema, basado en leyes físicas inmutables (la gravedad), entrega el recurso de forma constante y autónoma. 

El "Human-in-the-loop" es nuestra vasija de barro. Creemos que nuestra supervisión es la que mantiene el sistema funcionando, cuando en realidad, nuestra intervención es el síntoma de que aún no hemos construido nuestro acueducto. Salir del ciclo no es abandonar la responsabilidad; es elevar nuestra labor de la carga física a la ingeniería de sistemas.